\chapter{Evaluacija podsustava za \mbox{pohranu}}

Cilj ovog poglavlja jest usporediti dva ispitana podsustava za pohranu opisana u poglavljima arhitekture i implementacije. Vremenski orijentiranu bazu InfluxDB
s agentom Telegraf te dokumentno orijentiranu bazu Elasticsearch s agentom Filebeat. Usporedba se provodi nad istovjetnim ulaznim podatcima kako bi se
uočene razlike mogle pripisati isključivo pozadinskoj bazi i njezinu agentu, a ne ulazu. Naglasak je na trima mjerljivim svojstvima koja su za ciljanu primjenu
(potrošački laptop ograničenih resursa) presudna - zauzeću prostora na disku, potrošnji radne memorije te kašnjenju upita.

\section{Metodologija usporedbe}

Mjerenja su provedena na domaćinu temeljenom na jezgri \texttt{7.0.1 cachyos} (\texttt{x86\_64}) s ukupno približno 12 GB dostupne radne memorije.
Korištene su inačice InfluxDB 2.7, Elasticsearch 8.17 i Grafana 11.4, dok je veličina JVM gomile Elasticsearcha ograničena postavkom
\texttt{ES\_JAVA\_OPTS=-Xms512m} odnosno 512 MB.

Za razliku od produkcijske konfiguracije, u kojoj se zbog sukoba na portu 3000 (Grafana) i istovjetnog naziva senzorskog kontejnera izvodi isključivo jedan
podsustav, mjerenja su provedena uz istovremeno podignute obje baze i pripadne agente, no bez Grafane i bez senzora. Time se obje baze izlažu istom opterećenju
domaćina u istom trenutku, pa su mjerenja međusobno usporediva pod identičnim uvjetima. Kao ulaz korišten je jedan zamrznuti zapis \texttt{conn.log} veličine
4.7 MB s 14027 zapisa veza, koji oba agenta čitaju iz istoga dijeljenog direktorija.

Poštenost usporedbe počiva na tome da oba podsustava obrade isti ulaz, svaki ga pročitavši točno jednom, a to se provjerava brojem zapisa. Elasticsearch indeksira
svaki redak kao zaseban dokument, pa broj dokumenata mora odgovarati broju redaka ulaza. InfluxDB, naprotiv, točku jednoznačno određuje kombinacijom mjerenja,
skupa oznaka i vremenske oznake. Zapisi s istovjetnim vrijednostima tih polja stapaju se u jednu točku umjesto da se pišu zasebno. Posljedica je da InfluxDB
prijavljuje nešto manji broj točaka od broja ulaznih redaka, što nije gubitak podataka nego svojstvo modela pohrane i podrobnije se razmatra u raspravi.
Izmjereni su brojevi prikazani u tablici \ref{tab:verify}.

\begin{table}[ht]
    \centering
    \begin{tabular}{|c|c|}
        \hline
        \textbf{Izvor} & \textbf{Broj zapisa} \\
        \hline
        Redaka u \texttt{conn.log} (ulaz) & 14027 \\
        Elasticsearch dokumenata & 14027 \\
        InfluxDB točaka (\texttt{zeek\_conn}) & 13435 \\
        \hline
    \end{tabular}
    \caption{Provjera istovjetnosti ulaza}
    \label{tab:verify}
\end{table}

Mjerenja su automatizirana skriptom \texttt{benchmark.sh}, koja redom provodi provjeru ulaza, mjerenje prostora, uzorkovanje memorije i mjerenog kašnjenja upita.

\section{Potrošnja prostora na disku}

Mjeri se zauzeće na disku za istovjetan skup Zeekovih zapisa, uz izvornu datoteku \texttt{conn.log} kao referentnu vrijednost. Za InfluxDB mjeri se
veličina podatkovnog direktorija pripadnog spremnika \engl{bucket}, a za Elasticsearch veličina indeksa \texttt{zeek-conn} dobivena iz sučelja
\texttt{\_cat/indices}.

\begin{table}[h]
  \centering
  \begin{tabular}{|c|c|}
      \hline
      Izvor & Zauzeće \\
      \hline
      Sirovi \texttt{conn.log} referenca & 4.7 MB \\
      Elasticsearch indeks \texttt{zeek-conn} & 4.8 MB \\
      InfluxDB spremnik \texttt{zeek-logs} & 1.5 MB \\
      \hline
  \end{tabular}
  \caption{Zauzeće prostora na disku}
  \label{tab:storage}
\end{table}

Kao dodatni kontekst pri tumačenju ovih vrijednosti potrebno je naglasiti veličinu izmjerenog skupa. Drugim riječima izmjereni skup je malen (reda veličine nekoliko megabajta), 
pa razlike u zauzeću prostora prouzročavaju ponajprije dodatni troškovi formata pohrane, a ne ponašanje pri rastu količine podataka. Elasticsearch gradi invertirani indeks (uz pohranu podataka) 
koji radi pretrage po sadržaju polja, što povećava zauzeće diska u odnosu na sam volumen ulaza. InfluxDB, naprotiv, primjenjuje kompresiju prilagođenu vremenskim nizovima, pa je
izmjereni spremnik (1.5 MB) zauzeo otprilike trećinu Elasticsearchova indeksa i manje od same sirove datoteke. Posljedice tih značajki baza podrobnije se razmatraju u raspravi.

\section{Potrošnja radne memorije}

Potrošnja radne memorije mjerena je po kontejneru, u tri uzorka nakon što se sustav prilagodio, korištenjem Docker statistike. Tablica \ref{tab:ram}
prikazuje reprezentativne vrijednosti.

\begin{table}[h]
\centering
    \begin{tabular}{|c|c|}
    \hline
    \textbf{Kontejner} & \textbf{Radna memorija} \\
    \hline
    InfluxDB (baza, podsustav A) & 106 MB \\
    Telegraf (agent, podsustav A) & 40 MB \\
    Elasticsearch (baza, podsustav B) & 576 MB \\
    Filebeat (agent, podsustav B) & 84 MB \\
    \hline
    \end{tabular}
    \caption{Potrošnja radne memorije po kontejneru (nakon prilagodbe)}
    \label{tab:ram}
\end{table}

Razlika u potrošnji memorije najizraženija je i najpostojanija od svih mjerenih veličina. Cijeli podsustav s InfluxDB-om (baza i agent) zadržava se na razini
oko 146 MB, dok podsustav s Elasticsearchom (baza i agent) troši oko 660 MB, odnosno približno četiri i pol puta više. Bitno je da ta razlika nije posljedica
količine podataka, nego donje granice koju nameće Java virtualni stroj. Potrošnja Elasticsearcha određena je prije svega zadanom veličinom JVM gomile
(ovdje 512 MB), koja predstavlja donju granicu neovisnu o broju pohranjenih zapisa. Uz veću gomilu otisak bi bio i veći. Za ciljanu primjenu na sklopovlju
ograničenih resursa ta je donja granica presudna i razmatra se u zaključku. Valja napomenuti da je ta granica pomična: što se više memorije dodijeli, pretraživanja su brža.

\section{Kašnjenje upita}

Mjereno je kašnjenje triju upita koje koristi i nadzorna ploča: brojanje veza, ljestvica deset najčešćih odredišnih adresa te agregacija broja veza po
vremenskim prozorima. Svaki je upit pokrenut deset puta, uz odbacivanje dvaju početnih pokretanja radi zagrijavanja, a prijavljeni su medijan te najmanja i
najveća izmjerena vrijednost. Mjeri se ukupno kašnjenje kakvo klijent vidi na povratnoj adresi, dakle uključujući obradu zahtjeva, a ne samo unutarnje vrijeme
izvođenja upita u jezgri baze. Kako bi usporedba bila poštena, brojanje u oba sustava obavlja prolaz kroz zapise (u Elasticsearchu agregacijom
\texttt{value\_count}, a ne očitanjem održavanoga broja dokumenata).

\begin{table}[h]
    \centering
    \begin{tabular}{|c|c|c|}
        \hline
        \textbf{Upit} & \textbf{InfluxDB (Flux)} & \textbf{Elasticsearch} \\
        \hline
        Brojanje (sken) & 49,5 (45,3-51,9) & 4,4 (4,0-5,3) \\
        Deset najčešćih odredišta & 42,4 (40,3-44,5) & 5,1 (4,4-5,3) \\
        Agregacija po minutama & 63,6 (62,8-64,9) & 7,4 (6,1-9,1) \\
        \hline
    \end{tabular}
    \caption{Kašnjenje upita: medijan (najmanje-najveće), u milisekundama}
    \label{tab:latency}
\end{table}

Na svim trima upitima Elasticsearch postiže niže kašnjenje. Rezultat je nužno tumačiti uz činjenicu da je izmjereni skup malen, pa se mjeri pretežno dodatni trošak
izvođenja upita (priprema Flux upita i obrada zahtjeva nad sučeljem InfluxDB-a naspram Luceneova izvođenja u Elasticsearchu), a ne ponašanje pri rastu količine
podataka. Apsolutne su vrijednosti u oba sustava niske i za potrebe osvježavanja nadzorne ploče u sekundnim razmacima posve dostatne; razlika bi mogla dobiti na
značaju tek pri znatno većim skupovima ili pri upotrebi za pravovremenu reakciju, što ostaje smjernicom za daljnji rad.

\section{Rasprava rezultata}

Mjerenja se mogu objediniti kroz razliku u modelu podataka dvaju sustava, opisanu u poglavlju o teorijskim osnovama. Elasticsearch svaki Zeekov zapis
tretira kao dokument i nad poljima gradi invertirani indeks, što mu daje snažnu pretragu punim tekstom, povoljno kašnjenje agregacijskih upita i korelaciju
izvora, ali po cijeni većega prostornog otiska i, presudno, donje granice radne memorije koju nameće JVM. InfluxDB je optimiziran za vremenske
nizove pa mu je vrijeme glavna organizacijska dimenzija. Također primjenjuje kompresiju i politike zadržavanja, a memorijski otisak ostaje malen. Iz istoga modela slijedi 
i opažena razlika u broju zapisa iz tablice \ref{tab:verify}. Budući da InfluxDB točku određuje skupom oznaka i vremenskom oznakom, veze koje dijele te vrijednosti
stapaju se u jednu točku, pa je broj točaka manji od broja ulaznih redaka. Stapanje pogađa isključivo agregatni broj zapisa veza (otprilike 4\%) i posljedica je svjesne odluke da se 
identifikator veze \texttt{uid} drži poljem, a ne oznakom - držanje oznakom očuvalo bi svaki redak, ali bi količina vremenskih serija narasla na broj veza i poništila memorijsku prednost 
zbog koje se InfluxDB i odabire. Zapisi \texttt{notice} nisu zahvaćeni. Za primjenu usmjerenu na vizualizaciju obrazaca i dojavu, a ne na forenzičku rekonstrukciju pojedine veze, riječ je o 
prihvatljivom kompromisu. Unutar Elasticsearcha, gdje je svaki redak zaseban dokument, takvoga stapanja nema.

Posljedica je očita specijalizacija alata. InfluxDB je bolji izbor kada su upiti temeljeni na vremenu i kad su računalni resursi domaćina oskudni. Kad su
potrebni pretraga po sadržaju polja (primjerice traženje domena u \texttt{dns.log}), niže kašnjenje složenijih agregacija i korelacija više vrsta
zapisa u jednom upitu Elasticsearch se ističe kao očiti odabir. InfluxDB ne podržava prirodno te mogućnosti te zahtijeva agregaciju na strani klijenta.

Naposljetku valja istaknuti glavno ograničenje cijele evaluacije. Mjerenja su provedena na malenom skupu podataka, pa pouzdano govore o fiksnim zapisima i
donjim granicama potrošnje, ali ne i o ponašanju sustava pri znatnom rastu količine prometa. Mjerenje na većem skupu ostavlja se kao smjernica za daljnji
rad iako se može primijetiti trend gdje bi Elasticsearch bio puno bolja opcija.

\section{Zaključak usporedbe}

Za ciljanu primjenu ovog rada, nadzor lokalne mreže na potrošačkom sklopovlju ograničenih resursa, odlučujućom se pokazuje potrošnja radne memorije. Podsustav s InfluxDB-om zadržava se na 
razini stotinjak megabajta, dok podsustav s Elasticsearchom troši višestruko više, ponajprije zbog donje granice koju nameće JVM gomila. Ta je granica neovisna o količini pohranjenih podataka, 
pa Elasticsearch i pri zanemarivu opterećenju zauzima više memorije nego cijeli InfluxDB podsustav pod radnim opterećenjem; na domaćinu koji istovremeno obavlja uloge usmjernika, pristupne točke 
i senzora upravo je taj fiksni prag presudan. Tu je granicu moguće sniziti postavkom veličine JVM gomile, no na uštrb brzine i stabilnosti pretraživanja pri većim opterećenjima, čime slabi upravo prednost zbog koje bi se Elasticsearch inače odabrao.

Na svim trima izmjerenim upitima Elasticsearch postiže niže kašnjenje, kako prikazuje tablica \ref{tab:latency}. Apsolutne su vrijednosti obaju sustava 
ipak reda nekoliko desetaka milisekundi i za osvježavanje nadzorne ploče u sekundnim razmacima zanemarive, pa ne mijenjaju odluku koju daje memorijski otisak. Zauzeće diska na izmjerenom je 
skupu premaleno da bi samo po sebi bilo argumentom te se u ovom zaključku ne uzima kao odlučujuće.

Stoga se za produkcijsku konfiguraciju, u skladu s najavom iz poglavlja arhitekture da se odabire jedan od dvaju podsustava, bira InfluxDB. Time se ujedno pokazuje da je usporedivu 
temeljnu funkcionalnost moguće postići uz memorijski otisak reda veličine stotinjak megabajta, znatno ispod zahtjeva sveobuhvatnih rješenja poput Security Oniona navedenih u poglavlju o 
teorijskim osnovama.

Elasticsearch ostaje prikladan kad su potrebni pretraga po sadržaju polja (primjerice traženje domena u \texttt{dns.log}), niže kašnjenje složenijih agregacija i korelacija više vrsta 
zapisa u jednom upitu. Te prednosti slijede iz modela podataka, no nisu zasebno izmjerene u ovom radu jer mjerenja obuhvaćaju samo \texttt{conn.log} te se njihova kvantifikacija ostavlja 
kao smjernica za daljnji rad.